% figures03.tex -- Chapter 3 ("Concerning Lines") figures redrawn in TikZ.
%
% Fig 3-1 / 3-2 geometry is DIGITIZED from scan04.pdf p4 (book p.44): the six
% straight lines of Fig 3-1 were recovered by a Hough fit over the thresholded
% scan crop, and Fig 3-2 is the same six lines carrying the printed hand-wobble
% (measured perpendicular deviation 8-15 px at 170 dpi ~= 0.04-0.075 cm).
% Remaining figures are reconstructions at printed proportion; constrained
% points are placed by construction (never by eye) -- see tools/constraints/ch03.py.

% ============================================================ Fig 3-1, 3-2
% Six lines. L1 baseline carries A,B; L3 carries A',B'.
% AB and A'B' are drawn the SAME length in 3-1 (the text's whole point).
\newcommand{\FIGIIIONETWO}{\begin{bkfigure}[\linewidth]
\begin{minipage}{0.48\linewidth}\centering
\begin{tikzpicture}[scale=0.62]
  % digitized endpoints (scan px / 200, y flipped)
  \coordinate (p1) at (0.18,0.985);  \coordinate (q1) at (8.69,0.80);
  \coordinate (p2) at (0.905,3.195); \coordinate (q2) at (7.945,2.315);
  \coordinate (p3) at (4.06,4.075);  \coordinate (q3) at (7.60,0.81);
  \coordinate (p4) at (4.64,4.08);   \coordinate (q4) at (1.81,0.625);
  \coordinate (p5) at (2.005,3.565); \coordinate (q5) at (3.25,0.005);
  \coordinate (p6) at (7.35,3.64);   \coordinate (q6) at (5.74,0.765);
  \foreach \a/\b in {p1/q1,p2/q2,p3/q3,p4/q4,p5/q5,p6/q6} \draw[fig] (\a)--(\b);
  % A,B on the baseline; A',B' on L3 -- equal separations by construction
  \coordinate (A) at ($(p1)!0.085!(q1)$);  \coordinate (B) at ($(p1)!0.150!(q1)$);
  % A'B' drawn the SAME LENGTH as AB: t-step scaled by |L1|/|L3|.
  % Both slid down L3 past its crossing with L6 (t=0.658) so the labels have
  % clear air on the right-hand side of the line.
  \coordinate (Ap) at ($(p3)!0.83489!(q3)$); \coordinate (Bp) at ($(p3)!0.72000!(q3)$);
  \foreach \p in {A,B,Ap,Bp} {\dt{\p}}
  \node[vlab,above] at (A) {$A$};   \node[vlab,above] at (B) {$B$};
  \node[vlab,right,yshift=4pt] at (Ap) {$A'$};
  \node[vlab] at ($(Bp)+(47.3:0.45)$) {$B'$};
\end{tikzpicture}
\figcap{3--1}
\end{minipage}\hfill
\begin{minipage}{0.48\linewidth}\centering
\begin{tikzpicture}[scale=0.62]
  % same six lines seen "through uneven glass": wobble traced from the scan
  \draw[fig] plot[smooth,tension=0.85] coordinates
    {(0.245,0.745)(1.30,0.79)(2.35,0.755)(3.40,0.815)(4.45,0.80)(5.50,0.86)(6.45,0.90)(7.40,0.945)};
  \draw[fig] plot[smooth,tension=0.85] coordinates
    {(1.005,2.98)(1.90,2.90)(2.80,2.90)(3.70,2.79)(4.60,2.68)(5.50,2.62)(6.45,2.46)(7.39,2.355)};
  \draw[fig] plot[smooth,tension=0.85] coordinates
    {(4.05,3.87)(4.52,3.42)(4.95,2.98)(5.40,2.55)(5.85,2.09)(6.30,1.68)(6.75,1.30)(7.185,0.905)};
  \draw[fig] plot[smooth,tension=0.85] coordinates
    {(4.515,3.98)(4.14,3.53)(3.73,3.04)(3.38,2.55)(2.99,2.03)(2.63,1.53)(2.29,1.03)(1.945,0.515)};
  \draw[fig] plot[smooth,tension=0.85] coordinates
    {(1.97,3.355)(2.11,2.92)(2.28,2.49)(2.40,2.06)(2.56,1.64)(2.72,1.21)(2.93,0.79)(3.09,0.37)};
  \draw[fig] plot[smooth,tension=0.85] coordinates
    {(6.835,3.555)(6.62,3.11)(6.42,2.66)(6.24,2.20)(6.02,1.75)(5.83,1.30)(5.63,0.85)(5.43,0.40)};
  % same two pairs as Fig 3-1, at the same fractions along their curves
  \coordinate (A) at (0.83,0.775);  \coordinate (B) at (1.28,0.79);
  \coordinate (Ap) at (6.681,1.359); \coordinate (Bp) at (6.319,1.664);
  \foreach \p in {A,B,Ap,Bp} {\dt{\p}}
  \node[vlab,above] at (A) {$A$};   \node[vlab,above] at (B) {$B$};
  \node[vlab,right,yshift=4pt] at (Ap) {$A'$};
  \node[vlab] at ($(Bp)+(49.8:0.50)$) {$B'$};
\end{tikzpicture}
\figcap{3--2}
\end{minipage}
\end{bkfigure}}

% ======================================================= Fig 3-3, 3-4, 3-5
\newcommand{\FIGIIITHREEFOURFIVE}{\begin{bkfigure}[\linewidth]
\begin{minipage}{0.30\linewidth}\centering
\begin{tikzpicture}[scale=0.62]
  \draw[fig] (0,0) -- (5.0,0.55); \node[vlab,right] at (5.0,0.55) {$l$};
  \coordinate (P) at (2.0,0.22); \dt{P}
  \node[vlab,above] at (P) {$P$};
\end{tikzpicture}
\figcap{3--3}
\end{minipage}\hfill
\begin{minipage}{0.30\linewidth}\centering
\begin{tikzpicture}[scale=0.62]
  \coordinate (la) at (0,0.10); \coordinate (lb) at (4.8,1.42);
  \coordinate (ma) at (0,1.55); \coordinate (mb) at (4.8,0.02);
  \draw[fig] (la) -- (lb); \node[vlab,right] at (lb) {$l$};
  \draw[fig] (ma) -- (mb); \node[vlab,right] at (mb) {$m$};
  \coordinate (P) at (intersection of la--lb and ma--mb);
  \dt{P} \node[vlab,below] at (P) {$P$};
\end{tikzpicture}
\figcap{3--4}
\end{minipage}\hfill
\begin{minipage}{0.30\linewidth}\centering
\begin{tikzpicture}[scale=0.62]
  \coordinate (A) at (2.35,2.55); \coordinate (B) at (0.75,0.32);
  \coordinate (C) at (4.25,0.32);
  \draw[fig] ($(A)!-0.16!(B)$) -- ($(B)!-0.16!(A)$);
  \draw[fig] ($(A)!-0.16!(C)$) -- ($(C)!-0.16!(A)$);
  \draw[fig] ($(B)!-0.20!(C)$) -- ($(C)!-0.20!(B)$);
  % labels on the bisector of the free wedge at each vertex (the book puts
  % A above the crossing, B and C just inside the base)
  \node[vlab] at ($(A)+(92.5:0.75)$) {$A$};
  \node[vlab] at ($(B)+(27:0.90)$)   {$B$};
  \node[vlab] at ($(C)+(155:0.95)$)  {$C$};
\end{tikzpicture}
\figcap{3--5}
\end{minipage}
\end{bkfigure}}

% ============================================================ Fig 3-6, 3-7
% 3-6: two CURVED "lines" l and m meeting at A and B (the impossible picture).
\newcommand{\FIGIIISIXSEVEN}{\begin{bkfigure}[\linewidth]
\begin{minipage}{0.48\linewidth}\centering
\begin{tikzpicture}[scale=0.714]
  \coordinate (A) at (0.45,0.55); \coordinate (B) at (4.65,1.15);
  \draw[fig] (A) .. controls (1.7,1.55) and (3.3,1.75) .. (B);
  \draw[fig] (A) .. controls (1.7,0.10) and (3.3,0.35) .. (B);
  % BOTH curves run on past BOTH meeting points, so A and B each read as a
  % crossing (as in the book), not as a cusp
  \draw[fig] ($(A)!-0.09!(1.7,1.55)$) -- (A);
  \draw[fig] ($(A)!-0.09!(1.7,0.10)$) -- (A);
  \draw[fig] (B) -- ($(B)!-0.09!(3.3,1.75)$);
  \draw[fig] (B) -- ($(B)!-0.09!(3.3,0.35)$);
  \dt{A} \dt{B}
  \node[vlab,left] at (A) {$A$}; \node[vlab] at ($(B)+(100:0.50)$) {$B$};
  \node[vlab,above right] at (3.92,1.50) {$l$};
  \node[vlab,below right] at (4.55,0.88) {$m$};
\end{tikzpicture}
\figcap{3--6}
\end{minipage}\hfill
\begin{minipage}{0.48\linewidth}\centering
\begin{tikzpicture}[scale=0.714]
  \coordinate (C) at (2.55,2.60); \coordinate (A) at (0.55,0.35);
  \coordinate (B) at (4.35,0.62);
  \draw[fig] ($(A)!-0.18!(B)$) -- ($(B)!-0.16!(A)$);
  \draw[fig] ($(C)!-0.20!(A)$) -- ($(A)!-0.10!(C)$);
  \draw[fig] ($(C)!-0.20!(B)$) -- ($(B)!-0.10!(C)$);
  % m = line BC (extends up-LEFT above C); n = line CA (extends up-RIGHT).
  % C sits in the wedge ABOVE the crossing, between m and n, as in the book.
  \node[vlab] at ($(C)+(90:0.72)$) {$C$};
  \node[vlab,above left] at (A) {$A$}; \node[vlab,below right] at (B) {$B$};
  \node[vlab,left]  at ($(C)!-0.20!(B)$) {$m$};
  \node[vlab,right] at ($(C)!-0.20!(A)$) {$n$};
  \node[vlab,right] at ($(B)!-0.16!(A)$) {$l$};
\end{tikzpicture}
\figcap{3--7}
\end{minipage}
\end{bkfigure}}

% ============================================================ Fig 3-8, 3-9
\newcommand{\FIGIIIEIGHT}{\begin{bkfigure}[0.62\linewidth]
\begin{tikzpicture}[scale=0.806]
  \coordinate (P) at (0.40,0.35); \coordinate (Q) at (4.55,1.35);
  \draw[fig] (P) .. controls (1.7,1.35) and (3.2,1.70) .. (Q);
  \draw[fig] (P) .. controls (1.8,0.05) and (3.4,0.55) .. (Q);
  \draw[fig] ($(P)!-0.10!(1.7,1.35)$) -- (P);
  \draw[fig] ($(P)!-0.10!(1.8,0.05)$) -- (P);
  \draw[fig] (Q) -- ($(Q)!-0.10!(3.2,1.70)$);
  \draw[fig] (Q) -- ($(Q)!-0.10!(3.4,0.55)$);
  \dt{P} \dt{Q}
  \node[vlab,left] at (P) {$P$}; \node[vlab,above left] at (Q) {$Q$};
  \node[vlab,above right] at (4.30,1.60) {$a$};
  \node[vlab,below right] at (4.55,1.12) {$b$};
\end{tikzpicture}
\figcap{3--8}
\end{bkfigure}}

% 3-9: figure-eight -- two curved "lines" l and m crossing twice, A..E marked.
\newcommand{\FIGIIININE}{\begin{bkfigure}[0.62\linewidth]
\begin{tikzpicture}[scale=0.806]
  \coordinate (A) at (0.30,1.05); \coordinate (B) at (2.45,1.05);
  \coordinate (C) at (4.60,0.92);
  % D and E sit ON the right lobe: Bezier midpoints of the two arcs B..C
  \coordinate (D) at (3.5625,1.74625);
  \coordinate (E) at (3.5625,0.32125);
  \draw[fig] (A) .. controls (0.95,2.05) and (1.85,2.00) .. (B)
                 .. controls (3.05,2.05) and (4.10,1.95) .. (C);
  \draw[fig] (A) .. controls (0.95,0.05) and (1.85,0.10) .. (B)
                 .. controls (3.05,0.05) and (4.10,0.15) .. (C);
  \draw[fig] ($(A)!-0.12!(0.95,2.05)$) -- (A);
  \draw[fig] ($(A)!-0.12!(0.95,0.05)$) -- (A);
  \draw[fig] (C) -- ($(C)!-0.14!(4.10,1.95)$);
  \draw[fig] (C) -- ($(C)!-0.14!(4.10,0.15)$);
  \foreach \p in {A,B,C,D,E} {\dt{\p}}
  \node[vlab,above left] at (A) {$A$}; \node[vlab] at ($(B)+(90:0.58)$) {$B$};
  \node[vlab,below] at (C) {$C$};
  \node[vlab,above] at (D) {$D$}; \node[vlab,below] at (E) {$E$};
  \node[vlab,above right] at (4.80,1.36) {$l$};
  \node[vlab,right] at (5.02,0.62) {$m$};
\end{tikzpicture}
\figcap{3--9}
\end{bkfigure}}

% =============================================================== Fig 3-10
\newcommand{\FIGIIITEN}{\begin{bkfigure}[0.55\linewidth]
\begin{tikzpicture}[scale=0.806]
  \coordinate (A) at (2.45,2.85); \coordinate (B) at (0.35,0.30);
  \coordinate (C) at (4.55,0.55);
  \draw[fig] (A)--(B)--(C)--cycle;
  % D on AB, E on AC; the drawn line DE extended both ways
  \coordinate (D) at ($(A)!0.52!(B)$); \coordinate (E) at ($(A)!0.52!(C)$);
  \draw[fig] ($(D)!-0.30!(E)$) -- ($(E)!-0.22!(D)$);
  \foreach \p in {A,B,C,D,E} {\dt{\p}}
  \node[vlab,above] at (A) {$A$}; \node[vlab,below left] at (B) {$B$};
  \node[vlab,right] at (C) {$C$};
  \node[vlab,above left] at (D) {$D$}; \node[vlab,above right] at (E) {$E$};
\end{tikzpicture}
\figcap{3--10}
\end{bkfigure}}

% ========================================================= Fig 3-11, 3-12
\newcommand{\FIGIIIELEVENTWELVE}{\begin{bkfigure}[\linewidth]
\begin{minipage}{0.48\linewidth}\centering
\begin{tikzpicture}[scale=0.762]
  \draw[fig] (0,0) -- (4.6,0.62);
  \coordinate (A) at (0.85,0.115); \coordinate (B) at (2.95,0.398);
  \coordinate (C) at (3.55,0.478);
  \foreach \p in {A,B,C} {\dt{\p}}
  \node[vlab,above] at (A) {$A$}; \node[vlab,above] at (B) {$B$};
  \node[vlab,above] at (C) {$C$};
\end{tikzpicture}
\figcap{3--11}
\end{minipage}\hfill
\begin{minipage}{0.48\linewidth}\centering
\begin{tikzpicture}[scale=0.762]
  \draw[fig] (0,0) -- (4.6,0.62);
  \coordinate (A) at (0.60,0.081); \coordinate (D) at (2.05,0.276);
  \coordinate (B) at (2.60,0.350); \coordinate (C) at (4.05,0.546);
  \foreach \p in {A,D,B,C} {\dt{\p}}
  \node[vlab,above] at (A) {$A$}; \node[vlab,above] at (D) {$D$};
  \node[vlab,above] at (B) {$B$}; \node[vlab,above] at (C) {$C$};
\end{tikzpicture}
\figcap{3--12}
\end{minipage}
\end{bkfigure}}

% ========================================================= Fig 3-13, 3-14
\newcommand{\FIGIIITHIRTEENFOURTEEN}{\begin{bkfigure}[\linewidth]
\begin{minipage}{0.48\linewidth}\centering
\begin{tikzpicture}[scale=0.741]
  \draw[fig] (0,0) -- (4.6,1.55); \node[vlab,right] at (4.6,1.55) {$l$};
  \coordinate (P) at (0.55,1.95); \coordinate (Q) at (2.25,1.95);
  \draw[fig] (P) -- (Q);
  \dt{P} \dt{Q}
  \node[vlab,above] at (P) {$P$}; \node[vlab,above] at (Q) {$Q$};
  \node[slab,below] at ($(P)!0.5!(Q)$) {Same side};
\end{tikzpicture}
\figcap{3--13}
\end{minipage}\hfill
\begin{minipage}{0.48\linewidth}\centering
\begin{tikzpicture}[scale=0.741]
  \coordinate (P) at (1.25,2.10); \coordinate (Q) at (2.85,0.20);
  \coordinate (la) at (0.30,0.05); \coordinate (lb) at (4.30,2.30);
  \draw[fig] (la) -- (lb); \node[vlab,right] at (lb) {$l$};
  \draw[fig] (P) -- (Q);
  % R is where PQ meets l -- by construction, never by eye
  \coordinate (R) at (intersection of la--lb and P--Q);
  \foreach \p in {P,Q,R} {\dt{\p}}
  \node[vlab,above] at (P) {$P$}; \node[vlab,below] at (Q) {$Q$};
  \node[vlab] at ($(R)+(170:0.80)$) {$R$};
  \node[slab,right,align=left] at (3.05,1.14) {Opposite\\sides};
\end{tikzpicture}
\figcap{3--14}
\end{minipage}
\end{bkfigure}}

% =============================================================== Fig 3-15
% Halftone-style line illustration of three men standing in a line: PLATE.
\newcommand{\FIGIIIFIFTEEN}{\begin{bkfigure}[0.72\linewidth]
\setlength{\fboxsep}{10pt}%
\fbox{\begin{minipage}{0.86\linewidth}\centering\footnotesize
\textsc{Plate.}\ Illustration: three men standing in a line, marked
$A$, $B$, $C$, with the sight line from $A$ blocked at $B$.\\[2pt]
\emph{Original artwork -- not redrawn.}
\end{minipage}}
\figcap{3--15}
\end{bkfigure}}

% ========================================================= Fig 3-16, 3-17
\newcommand{\FIGIIISIXTEENSEVENTEEN}{\begin{bkfigure}[\linewidth]
\begin{minipage}{0.48\linewidth}\centering
\begin{tikzpicture}[scale=0.768]
  \coordinate (A) at (0.20,0.15); \coordinate (B) at (3.55,2.20);
  \draw[fig] (A) -- (B);
  \coordinate (P) at ($(A)!0.46!(B)$);
  \foreach \p in {A,B,P} {\dt{\p}}
  \node[vlab,below left] at (A) {$A$}; \node[vlab,above] at (B) {$B$};
  \node[vlab,above left] at (P) {$P$};
\end{tikzpicture}
\figcap{3--16}
\end{minipage}\hfill
\begin{minipage}{0.48\linewidth}\centering
\begin{tikzpicture}[scale=0.768]
  \coordinate (A) at (2.05,2.45); \coordinate (B) at (0.25,0.20);
  \coordinate (C) at (3.75,0.05);
  \draw[fig] (A)--(B)--(C)--cycle;
  \node[vlab,above] at (A) {$A$}; \node[vlab,left] at (B) {$B$};
  \node[vlab,below right] at (C) {$C$};
\end{tikzpicture}
\figcap{3--17}
\end{minipage}
\end{bkfigure}}

% =============================================================== Fig 3-18
% l meets side AB and misses BC; C is the apex. l passes through no vertex.
\newcommand{\FIGIIIEIGHTEEN}{\begin{bkfigure}[0.55\linewidth]
\begin{tikzpicture}[scale=0.806]
  \coordinate (C) at (1.85,2.60); \coordinate (B) at (0.20,0.30);
  \coordinate (A) at (4.30,0.35);
  \draw[fig] (C)--(B)--(A)--cycle;
  % l: enters through AB (the base) and leaves through AC
  \coordinate (X) at ($(B)!0.42!(A)$);      % on AB (base)
  \coordinate (Y) at ($(C)!0.52!(A)$);      % on AC
  % l runs from just below the base up past AC to about C's height, as printed
  \coordinate (ltop) at ($(Y)!-1.05!(X)$);
  \draw[fig] ($(X)!-0.35!(Y)$) -- (ltop);
  \node[vlab,above] at (C) {$C$}; \node[vlab,left] at (B) {$B$};
  \node[vlab,right] at (A) {$A$};
  \node[vlab,above right] at (ltop) {$l$};
\end{tikzpicture}
\figcap{3--18}
\end{bkfigure}}

% =============================================================== Fig 3-19
\newcommand{\FIGIIININETEEN}{\begin{bkfigure}[0.62\linewidth]
\begin{tikzpicture}[scale=0.806]
  \draw[fig] (0,0) -- (5.4,0); \node[vlab,right] at (5.4,0) {$l$};
  \coordinate (O) at (2.55,0); \dt{O}
  \node[vlab,below] at (O) {$O$};
  \draw[fig,line width=0.5pt] (0.35,0.30) -- (0.35,0.46) -- (2.45,0.46) -- (2.45,0.30);
  \draw[fig,line width=0.5pt] (2.65,0.30) -- (2.65,0.46) -- (4.95,0.46) -- (4.95,0.30);
  \node[slab,above] at (1.40,0.46) {One side};
  \node[slab,above] at (3.80,0.46) {Other side};
\end{tikzpicture}
\figcap{3--19}
\end{bkfigure}}

% ========================================================= Fig 3-20, 3-21
\newcommand{\FIGIIITWENTYTWENTYONE}{\begin{bkfigure}[\linewidth]
\begin{minipage}{0.48\linewidth}\centering
\begin{tikzpicture}[scale=0.751]
  \coordinate (la) at (0,0.55);     \coordinate (lb) at (4.35,0.55);
  \coordinate (ma) at (1.05,-0.45); \coordinate (mb) at (2.95,2.40);
  \draw[fig] (la) -- (lb); \node[vlab,right] at (lb) {$l$};
  \draw[fig] (ma) -- (mb); \node[vlab,above] at (mb) {$m$};
  % O is the point l and m have in common; A is a point of m, not of l
  \coordinate (O) at (intersection of la--lb and ma--mb);
  \coordinate (A) at ($(ma)!0.79!(mb)$);
  \dt{O} \dt{A}
  \node[vlab] at ($(O)+(298:0.55)$) {$O$}; \node[vlab,right] at (A) {$A$};
\end{tikzpicture}
\figcap{3--20}
\end{minipage}\hfill
\begin{minipage}{0.48\linewidth}\centering
\begin{tikzpicture}[scale=0.751]
  \coordinate (la) at (0,0.55);     \coordinate (lb) at (4.35,0.55);
  \coordinate (ma) at (1.05,-0.45); \coordinate (mb) at (2.95,2.40);
  \draw[fig] (la) -- (lb); \node[vlab,right] at (lb) {$l$};
  \draw[fig] (ma) -- (mb); \node[vlab,above] at (mb) {$m$};
  \coordinate (O) at (intersection of la--lb and ma--mb);
  \coordinate (A) at ($(ma)!0.79!(mb)$);
  % P and Q are points of l on opposite sides of O:  POQ
  \coordinate (Q) at ($(la)!0.138!(lb)$); \coordinate (P) at ($(la)!0.655!(lb)$);
  \foreach \p in {O,A,P,Q} {\dt{\p}}
  \node[vlab] at ($(O)+(298:0.55)$) {$O$}; \node[vlab,right] at (A) {$A$};
  \node[vlab,below] at (Q) {$Q$}; \node[vlab,below] at (P) {$P$};
\end{tikzpicture}
\figcap{3--21}
\end{minipage}
\end{bkfigure}}

% ========================================================= Fig 3-22, 3-23
\newcommand{\FIGIIITWENTYTWOTWENTYTHREE}{\begin{bkfigure}[\linewidth]
\begin{minipage}{0.48\linewidth}\centering
\begin{tikzpicture}[scale=0.755]
  \coordinate (O) at (1.05,0.35); \coordinate (lend) at (4.35,1.32);
  % the part of l that is NOT in the ray, dashed: the exact opposite ray
  \draw[fig,dashed,dash pattern=on 3pt off 2.4pt] ($(O)!-0.318!(lend)$) -- (O);
  \draw[fig] (O) -- (lend); \node[vlab,right] at (lend) {$l$};
  \coordinate (P) at ($(O)!0.394!(lend)$);
  \dt{O} \dt{P}
  \node[vlab,above] at (O) {$O$}; \node[vlab,above] at (P) {$P$};
  \node[slab,below] at (1.90,0.52) {Ray};
\end{tikzpicture}
\figcap{3--22}
\end{minipage}\hfill
\begin{minipage}{0.48\linewidth}\centering
\begin{tikzpicture}[scale=0.755]
  \coordinate (P) at (2.05,0.30);
  \draw[fig] (0.15,0.30) -- (4.35,0.30); \node[vlab,right] at (4.35,0.30) {$l$};
  \draw[fig] (P) -- (3.35,2.55); \node[vlab,above] at (3.35,2.55) {$r$};
  \dt{P} \node[vlab,above left] at (P) {$P$};
\end{tikzpicture}
\figcap{3--23}
\end{minipage}
\end{bkfigure}}

% ========================================================= Fig 3-24, 3-25
\newcommand{\FIGIIITWENTYFOURTWENTYFIVE}{\begin{bkfigure}[\linewidth]
\begin{minipage}{0.52\linewidth}\centering
\begin{tikzpicture}[scale=0.780]
  % (i) angle at O, upper left
  \coordinate (O1) at (0.30,2.05);
  \draw[fig] (O1) -- (1.85,2.95); \draw[fig] (O1) -- (1.85,2.05);
  \dt{O1} \node[vlab,left] at (O1) {$O$};
  % (ii) angle at O, lower left
  \coordinate (O2) at (0.62,1.35);
  \draw[fig] (O2) -- (2.05,1.05); \draw[fig] (O2) -- (-0.15,0.20);
  \dt{O2} \node[vlab,left] at (O2) {$O$};
  % (iii) angle AOB, right
  \coordinate (O3) at (3.55,2.20); \coordinate (A3) at (2.55,3.10);
  \coordinate (B3) at (2.85,1.25);
  \draw[fig] (O3) -- (A3); \draw[fig] (O3) -- (B3);
  \dt{O3} \dt{A3} \dt{B3}
  \node[vlab,right] at (O3) {$O$}; \node[vlab,above] at (A3) {$A$};
  \node[vlab,below] at (B3) {$B$};
  % (iv) straight angle
  \draw[fig] (0.85,0.00) -- (3.35,0.00);
  \coordinate (O4) at (2.05,0.00); \dt{O4}
  \node[vlab,below] at (O4) {$O$};
\end{tikzpicture}
\figcap{3--24}
\end{minipage}\hfill
\begin{minipage}{0.44\linewidth}\centering
\begin{tikzpicture}[scale=0.780]
  % (a) the arrowheads are on the ENDS OF THE ARC, pointing at the two sides;
  % the sides themselves carry no arrows.  Verified on scans/scan04.pdf p18
  % (book p.54), where both arrowheads are plainly on the arc.
  \coordinate (Oa) at (0.15,2.55);
  \draw[fig] (Oa) -- (1.90,3.15);
  \draw[fig] (Oa) -- (1.75,2.10);
  % every arc is centred on its vertex and stops exactly on the two sides
  \draw[fig,line width=0.5pt,<->,>=stealth]
    ($(Oa)+(18.93:1.05)$) arc (18.93:-15.71:1.05);
  \node[slab] at ($(Oa)+(1.61:1.30)$) {1};
  \node[slab,right] at (2.25,2.55) {(a)};
  % (b) plain sides
  \coordinate (Ob) at (0.15,1.35);
  \draw[fig] (Ob) -- (1.90,1.95); \draw[fig] (Ob) -- (1.75,0.90);
  \draw[fig,line width=0.5pt] ($(Ob)+(18.93:1.05)$) arc (18.93:-15.71:1.05);
  \node[slab] at ($(Ob)+(1.61:1.30)$) {1};
  \node[slab,right] at (2.25,1.35) {(b)};
  % (c) obtuse
  \coordinate (Oc) at (0.35,0.10);
  \draw[fig] (Oc) -- (2.10,0.10); \draw[fig] (Oc) -- (-0.25,0.95);
  \draw[fig,line width=0.5pt] ($(Oc)+(0:0.90)$) arc (0:125.18:0.90);
  \node[slab] at ($(Oc)+(62.59:0.58)$) {2};
  \node[slab,right] at (2.25,0.10) {(c)};
\end{tikzpicture}
\figcap{3--25}
\end{minipage}
\end{bkfigure}}

% =============================================================== Fig 3-26
% Two lines l, m crossing at O; four rays l1,l2,m1,m2; angles 1,2,3,4.
\newcommand{\FIGIIITWENTYSIX}{\begin{bkfigure}[0.66\linewidth]
\begin{tikzpicture}[scale=0.910]
  \coordinate (L1) at (4.85,2.05); \coordinate (L2) at (0.35,0.15);
  \coordinate (M1) at (4.85,0.35); \coordinate (M2) at (0.35,1.95);
  \draw[fig] (L2) -- (L1); \draw[fig] (M2) -- (M1);
  \coordinate (O) at (intersection of L2--L1 and M2--M1);
  \dt{O}
  \node[vlab,right] at (L1) {$l_1$}; \node[vlab,left]  at (L2) {$l$};
  \node[vlab,right] at (M1) {$m_1$}; \node[vlab,left]  at (M2) {$m$};
  % ray labels offset perpendicular to their own line, so they sit beside it
  \node[vlab] at ($($(M2)!0.28!(M1)$)+(70.42:0.42)$) {$m_2$};
  \node[vlab] at ($($(L2)!0.30!(L1)$)+(-67.10:0.42)$) {$l_2$};
  % the four arcs are centred on O and stop exactly on the four rays
  \draw[fig,line width=0.5pt] ($(O)+(160.44:0.62)$) arc (160.44:202.89:0.62);
  \draw[fig,line width=0.5pt] ($(O)+(-19.58:0.62)$) arc (-19.58:22.90:0.62);
  \draw[fig,line width=0.5pt] ($(O)+(22.90:0.62)$)  arc (22.90:160.44:0.62);
  \draw[fig,line width=0.5pt] ($(O)+(202.89:0.62)$) arc (202.89:340.42:0.62);
  % angle numerals on each region's bisector, clear of the arcs
  \node[slab] at ($(O)+(181.67:0.95)$) {1};
  \node[slab] at ($(O)+(1.66:0.95)$)   {2};
  \node[slab] at ($(O)+(91.67:0.95)$)  {3};
  \node[slab] at ($(O)+(271.66:0.95)$) {4};
  % the book sets O just below-left of the crossing, inside the ring of arcs;
  % at this type size the label padding will not clear the two lines there, so
  % it stays outside the ring, below-right (see ch03-UNCERTAIN.md)
  \node[vlab] at ($(O)+(303:1.00)$) {$O$};
\end{tikzpicture}
\figcap{3--26}
\end{bkfigure}}

% =============================================================== Fig 3-27
% Three drawings of a pair of angles sharing a vertex and one side.
\newcommand{\FIGIIITWENTYSEVEN}{\begin{bkfigure}[0.82\linewidth]
\begin{tikzpicture}[scale=0.780]
  % (i) narrow pair
  % the narrowest of the three pairs; opened from 18/24 deg to 23/25 so the
  % numerals clear both rays at this type size (the book's own fan is tight)
  \coordinate (O) at (0.15,0.55);
  \draw[fig] (O) -- (1.891,1.724); \draw[fig] (O) -- (2.261,0.960);
  \draw[fig] (O) -- (2.091,0.066);
  \draw[fig,line width=0.5pt] ($(O)+(34:1.25)$) arc (34:11:1.25);
  \draw[fig,line width=0.5pt] ($(O)+(11:1.25)$) arc (11:-14:1.25);
  \node[slab] at ($(O)+(22.5:1.62)$) {1};
  \node[slab] at ($(O)+(-1.5:1.62)$) {2};
\end{tikzpicture}\hspace{1.6em}
\begin{tikzpicture}[scale=0.780]
  % (ii) the book's middle drawing: ONE ray straight up, TWO rays down-left and
  % down-right (an inverted Y), the vertical being the common side.  The two
  % arcs bulge out either side of it -- measured off the printed figure.
  \coordinate (O) at (1.15,1.35);
  \draw[fig] (O) -- (1.15,2.75);            % common side, straight up
  \draw[fig] (O) -- (0.30,0.25);            % noncommon side, down-left
  \draw[fig] (O) -- (2.05,0.25);            % noncommon side, down-right
  \draw[fig,line width=0.5pt] ($(O)+(232.30:0.80)$) arc (232.30:90:0.80);
  \draw[fig,line width=0.5pt] ($(O)+(90:0.80)$) arc (90:-50.71:0.80);
  \node[slab] at ($(O)+(161.15:1.14)$) {1};
  \node[slab] at ($(O)+(19.65:1.14)$)  {2};
\end{tikzpicture}\hspace{1.6em}
\begin{tikzpicture}[scale=0.780]
  % (iii) wide pair over a base line
  \coordinate (O) at (1.25,0.25);
  \draw[fig] (O) -- (0.05,1.35); \draw[fig] (O) -- (1.55,2.05);
  \draw[fig] (O) -- (2.85,0.90);
  \draw[fig,line width=0.5pt] ($(O)+(137.48:1.35)$) arc (137.48:80.54:1.35);
  \draw[fig,line width=0.5pt] ($(O)+(80.54:1.35)$) arc (80.54:22.11:1.35);
  \node[slab] at ($(O)+(109.01:1.80)$) {1};
  \node[slab] at ($(O)+(51.33:1.80)$)  {2};
\end{tikzpicture}
\figcap{3--27}
\end{bkfigure}}

% =============================================================== Fig 3-28
\newcommand{\FIGIIITWENTYEIGHT}{\begin{bkfigure}[0.55\linewidth]
\begin{tikzpicture}[scale=0.858]
  \coordinate (O) at (2.05,0.20);
  \draw[fig] (0.20,0.20) -- (4.15,0.20);
  \draw[fig] (O) -- (3.55,1.70);          % common side, at 45 degrees
  % both arcs run from a side of the angle to the common side, no gap
  \draw[fig,line width=0.5pt] ($(O)+(180:1.20)$) arc (180:45:1.20);
  \node[slab] at ($(O)+(112.5:0.85)$) {1};
  \draw[fig,line width=0.5pt] ($(O)+(0:0.67)$) arc (0:45:0.67);
  \node[slab] at ($(O)+(22.5:1.00)$) {2};
\end{tikzpicture}
\figcap{3--28}
\end{bkfigure}}

% ========================================================= Fig 3-29, 3-30
% 3-29: lines l and m crossing at O; the four regions cut by them, two of
% which are hatched (vertical / horizontal) so their overlap reads doubly hatched.
\newcommand{\FIGIIITWENTYNINETHIRTY}{\begin{bkfigure}[\linewidth]
\begin{minipage}{0.48\linewidth}\centering
\begin{tikzpicture}[scale=0.730]
  % The book's picture: A B C D at the corners of a rectangle, l = AC and
  % m = BD crossing at O.  VERTICAL rules shade the side of l that contains D;
  % HORIZONTAL rules shade the side of m that contains C.  The two families
  % therefore overlap in triangle OCD, and that doubly-shaded wedge IS the
  % interior of angle COD -- which is the whole point of the figure.
  \coordinate (B) at (0.35,2.55); \coordinate (Cc) at (3.95,2.55);
  \coordinate (A) at (0.35,0.35); \coordinate (D) at (3.95,0.35);
  \coordinate (O) at (intersection of A--Cc and B--D);
  % Each rule is drawn at its exact length rather than clipped: a clipped rule
  % still carries its full path in the PDF, so the collision audit would see
  % ink where the picture shows none.
  % D's side of l = triangle A D C: vertical rules from AD up to line AC
  \foreach \x in {0.45,0.55,...,3.90}
    \draw[fig,line width=0.35pt] (\x,0.35) -- (\x,{0.35+(\x-0.35)*0.611111});
  % C's side of m = triangle B C D: horizontal rules from line BD across to DC
  \foreach \y in {0.45,0.55,...,2.50}
    \draw[fig,line width=0.35pt] ({0.35+(2.55-\y)*1.636364},\y) -- (3.95,\y);
  % the two lines: carried past the far corner to take the label, and past the
  % near point as well -- in the book A and B are POINTS ON the lines, not ends
  \draw[fig] ($(A)!-0.07!(Cc)$) -- ($(A)!1.10!(Cc)$);
  \node[vlab,right] at ($(A)!1.10!(Cc)$) {$l$};
  \draw[fig] ($(B)!-0.07!(D)$) -- ($(B)!1.10!(D)$);
  \node[vlab,right] at ($(B)!1.10!(D)$) {$m$};
  \foreach \p in {A,B,Cc,D,O} {\dt{\p}}
  % A and B sit clear of the tails the two lines now carry past them
  \node[vlab] at ($(A)+(150:0.50)$) {$A$};  \node[vlab] at ($(B)+(95:0.58)$) {$B$};
  \node[vlab,above left] at (Cc) {$C$}; \node[vlab,below left] at (D) {$D$};
  \node[vlab] at ($(O)+(180:0.85)$) {$O$};   % the one unshaded wedge, as in the book
\end{tikzpicture}
\figcap{3--29}
\end{minipage}\hfill
\begin{minipage}{0.48\linewidth}\centering
\begin{tikzpicture}[scale=0.730]
  % r_1 lies on l_1 and r_2 on l_2, so each dashed piece is the ray OPPOSITE
  % its solid one: built as a reflection through O, never placed by eye.
  \coordinate (O) at (2.30,1.35);
  \coordinate (R1) at (4.30,2.55); \coordinate (R2) at (4.10,0.10);
  \coordinate (L1) at ($(O)!-0.90!(R1)$);   % opposite ray of r_1
  \coordinate (L2) at ($(O)!-1.00!(R2)$);   % opposite ray of r_2
  \draw[fig] (O) -- (R1); \draw[fig] (O) -- (R2);
  \draw[aux] (O) -- (L2); \draw[aux] (O) -- (L1);
  \node[vlab,right] at (R1) {$r_1$}; \node[vlab,right] at (R2) {$r_2$};
  \node[vlab,left] at (L2) {$l_2$};  \node[vlab,left] at (L1) {$l_1$};
  \dt{O} \node[vlab] at ($(O)+(90:0.62)$) {$O$};
\end{tikzpicture}
\figcap{3--30}
\end{minipage}
\end{bkfigure}}

% =============================================================== Fig 3-31
\newcommand{\FIGIIITHIRTYONE}{\begin{bkfigure}[0.70\linewidth]
\begin{tikzpicture}[scale=0.806]
  \coordinate (O) at (2.05,1.35);
  \coordinate (R1) at (4.05,2.45); \coordinate (R2) at (3.95,0.15);
  \coordinate (L1) at ($(O)!-0.85!(R1)$);   % opposite ray of r_1
  \coordinate (L2) at ($(O)!-0.95!(R2)$);   % opposite ray of r_2
  \draw[fig] (O) -- (R1); \node[vlab,right] at (R1) {$r_1$};
  \draw[fig] (O) -- (R2); \node[vlab,right] at (R2) {$r_2$};
  % The book draws l_1 and l_2 SOLID here, unlike Fig 3-30 where the rays that
  % are not sides of the angle are dashed -- checked on scans/scan04.pdf p20.
  \draw[fig] (O) -- (L2); \node[vlab,left] at (L2) {$l_2$};
  \draw[fig] (O) -- (L1); \node[vlab,left] at (L1) {$l_1$};
  \dt{O} \node[vlab] at ($(O)+(90:0.62)$) {$O$};
  % the separate segment PQ with A between; it runs a little past Q
  \coordinate (P) at (5.15,0.60); \coordinate (Q) at (7.15,0.60);
  \draw[fig] (P) -- ($(P)!1.14!(Q)$);
  \coordinate (Aa) at ($(P)!0.5!(Q)$);
  \foreach \p in {P,Q,Aa} {\dt{\p}}
  \node[vlab,above] at (P) {$P$}; \node[vlab,above] at (Q) {$Q$};
  \node[vlab,above] at (Aa) {$A$};
\end{tikzpicture}
\figcap{3--31}
\end{bkfigure}}

% =============================================================== Fig 3-32
\newcommand{\FIGIIITHIRTYTWO}{\begin{bkfigure}[0.60\linewidth]
\begin{tikzpicture}[scale=0.806]
  \node[vlab,left] at (0.20,2.45) {$r_2'$};
  \node[vlab,left] at (0.25,0.20) {$r_1'$};
  % the two full lines through O; r1' / r2' are the rays opposite r1 / r2
  \coordinate (ta) at (0.20,2.45); \coordinate (tb) at (4.05,0.30);
  \coordinate (ua) at (0.25,0.20); \coordinate (ub) at (4.15,2.35);
  \draw[fig] (ta) -- (tb);
  \draw[fig] (ua) -- (ub);
  \coordinate (O) at (intersection of ta--tb and ua--ub);
  \draw[fig] (O) -- (4.20,1.34); \node[vlab,right] at (4.20,1.34) {$r$};
  \node[vlab,right] at (4.15,2.35) {$r_1$}; \node[vlab,right] at (4.05,0.30) {$r_2$};
  \dt{O} \node[vlab] at ($(O)+(90:0.52)$) {$O$};
  % the two arcs mark angle(r_1,r) and angle(r,r_2): centred on O, ending on
  % the rays (O = (2.2524,1.3039); r_1 at 28.88, r at 1.06, r_2 at -29.18 deg)
  \draw[fig,line width=0.5pt] ($(O)+(28.88:1.40)$) arc (28.88:1.06:1.40);
  \draw[fig,line width=0.5pt] ($(O)+(1.06:1.40)$)  arc (1.06:-29.18:1.40);
\end{tikzpicture}
\figcap{3--32}
\end{bkfigure}}

% =============================================================== Fig 3-33
\newcommand{\FIGIIITHIRTYTHREE}{\begin{bkfigure}[0.72\linewidth]
\begin{tikzpicture}[scale=0.806]
  % l2 : the shallow line, carrying r2 to the right and r2' to the left
  \coordinate (L2R) at (5.75,1.95); \coordinate (L2L) at (0.30,1.75);
  \draw[fig] (L2L) -- (L2R);
  \node[vlab,right] at (L2R) {$l_2$}; \node[vlab,left] at (L2L) {$r_2'$};
  % l1 : the steep line; O is where they meet
  \coordinate (L1T) at (4.15,3.35); \coordinate (L1B) at (1.95,0.15);
  \draw[fig] (L1B) -- (L1T);
  \coordinate (O) at (intersection of L1B--L1T and L2L--L2R);
  \node[vlab,right] at (L1T) {$l_1$};
  \draw[aux] (L1B) -- (1.70,-0.20); \node[vlab,below] at (1.70,-0.20) {$r_1'$};
  \node[vlab,left]  at ($(O)!0.72!(L1T)$) {$r_1$};   % ray r1: above O
  \node[vlab,below] at ($(O)!0.80!(L2R)$) {$r_2$};   % ray r2: right of O
  % P lies in the interior of the angle; Q_3 is across BOTH lines from it, so
  % the book's segment Q_3 P has to meet l_1 in a point R and l_2 in a point S.
  % R and S are taken as intersections -- the proof's whole content.
  \coordinate (P) at (5.15,2.45);  \coordinate (Qc) at (1.85,1.15);
  \draw[fig] (Qc) -- (P);
  \coordinate (R) at (intersection of Qc--P and L1B--L1T);
  \coordinate (S) at (intersection of Qc--P and L2L--L2R);
  \foreach \p in {O,P,S,R} {\dt{\p}}
  \node[vlab] at ($(O)+(118.8:0.55)$) {$O$};
  \node[vlab,above right] at (P) {$P$};
  \node[vlab] at ($(S)+(60:0.75)$) {$S$};
  \node[vlab] at ($(R)+(308.5:0.48)$) {$R$};
  % the five sample positions.  Q_1, Q_3, Q_5 lie in the three regions the
  % lines cut off; Q_2 lies ON the ray r_2' and Q_4 ON the ray r_1' -- that
  % distinction is what Exercises *12-*14 turn on.
  \coordinate (Qa) at (2.55,3.05);                    % Q_1
  \coordinate (Qb) at ($(L2L)!0.22!(L2R)$);           % Q_2, on r_2'
  \coordinate (Qd) at ($(L1B)!0.12!(L1T)$);           % Q_4, on r_1'
  \coordinate (Qe) at (4.30,0.80);                    % Q_5
  \foreach \p in {Qa,Qb,Qc,Qd,Qe} {\dt{\p}}
  \node[vlab,above] at (Qa) {$Q_1$}; \node[vlab,above] at (Qb) {$Q_2$};
  \node[vlab] at ($(Qc)+(235:0.55)$) {$Q_3$};
  \node[vlab,right] at (Qd) {$Q_4$}; \node[vlab,above] at (Qe) {$Q_5$};
\end{tikzpicture}
\figcap{3--33}
\end{bkfigure}}

% ========================================================= Fig 3-34, 3-35
\newcommand{\FIGIIITHIRTYFOURTHIRTYFIVE}{\begin{bkfigure}[\linewidth]
\begin{minipage}{0.48\linewidth}\centering
\begin{tikzpicture}[scale=0.749]
  \coordinate (B) at (1.75,2.35); \coordinate (A) at (0.20,0.25);
  \coordinate (C) at (3.75,0.25);
  \draw[fig] (A)--(B)--(C)--cycle;
  \node[vlab,above] at (B) {$B$}; \node[vlab,left] at (A) {$A$};
  \node[vlab,right] at (C) {$C$};
\end{tikzpicture}
\figcap{3--34}
\end{minipage}\hfill
\begin{minipage}{0.48\linewidth}\centering
\begin{tikzpicture}[scale=0.749]
  \coordinate (B) at (1.55,2.35); \coordinate (A) at (0.20,0.25);
  \coordinate (C) at (3.55,0.25);
  \draw[fig] (A)--(B)--(C)--cycle;
  \coordinate (P) at (1.45,1.05); \coordinate (Q) at (3.55,2.20);
  \draw[fig] (P) -- (Q);
  \dt{P} \dt{Q}
  \node[vlab,above] at (B) {$B$}; \node[vlab,left] at (A) {$A$};
  \node[vlab,right] at (C) {$C$};
  \node[vlab,above,yshift=2pt] at (P) {$P$}; \node[vlab,above right] at (Q) {$Q$};
\end{tikzpicture}
\figcap{3--35}
\end{minipage}
\end{bkfigure}}

% =============================================================== Fig 3-36
\newcommand{\FIGIIITHIRTYSIX}{\begin{bkfigure}[0.50\linewidth]
\begin{tikzpicture}[scale=0.806]
  \coordinate (O) at (0.75,0.35); \coordinate (A) at (3.85,0.15);
  \coordinate (Ot) at (2.55,2.85); \coordinate (At) at (1.55,2.75);
  \draw[fig,dashed,dash pattern=on 3pt off 2.4pt] ($(O)!-0.23!(A)$) -- (O);
  \draw[fig] (O) -- (A);
  \draw[fig] (O) -- (Ot);
  \draw[fig] (A) -- (At);
  % B is where the ray from O meets the other side of the angle at A
  \coordinate (B) at (intersection of O--Ot and A--At);
  \dt{O} \dt{A} \dt{B}
  \node[vlab,above left] at (O) {$O$}; \node[vlab,right] at (A) {$A$};
  \node[vlab] at ($(B)+(93:0.80)$) {$B$};
\end{tikzpicture}
\figcap{3--36}
\end{bkfigure}}

% =============================================================== Fig 3-37
\newcommand{\FIGIIITHIRTYSEVEN}{\begin{bkfigure}[0.52\linewidth]
\begin{tikzpicture}[scale=0.806]
  \coordinate (V) at (1.35,0.65);
  \coordinate (M1) at (3.35,2.75);
  \draw[fig] (V) -- (3.95,0.65); \node[vlab,right] at (3.95,0.65) {$l_1$};
  \draw[fig] (V) -- (M1); \node[vlab,above] at (M1) {$m_1$};
  % l and m are the REST of the two lines: exact opposite rays through V
  \draw[fig,dashed,dash pattern=on 3pt off 2.4pt] (V) -- (0.15,0.65);
  \node[vlab,left] at (0.15,0.65) {$l$};
  \coordinate (Mo) at ($(V)!-0.45!(M1)$);
  \draw[fig,dashed,dash pattern=on 3pt off 2.4pt] (V) -- (Mo);
  \node[vlab,below] at (Mo) {$m$};
  \draw[fig,line width=0.5pt] ($(V)+(0:1.10)$) arc (0:46.40:1.10);
\end{tikzpicture}
\figcap{3--37}
\end{bkfigure}}

% ========================================================= Fig 3-38, 3-39
\newcommand{\FIGIIITHIRTYEIGHTTHIRTYNINE}{\begin{bkfigure}[\linewidth]
\begin{minipage}{0.48\linewidth}\centering
\begin{tikzpicture}[scale=0.675]
  % sides of the angle solid, the opposite rays r_1' and r_2' dashed
  \coordinate (l2L) at (0.30,1.05); \coordinate (l2R) at (4.55,1.25);
  \coordinate (ka) at (1.65,0.15);  \coordinate (kb) at (3.35,2.85);
  \coordinate (O) at (intersection of l2L--l2R and ka--kb);
  \draw[fig] (O) -- (l2R); \node[vlab,right] at (l2R) {$l_2$};
  \draw[aux] (O) -- (l2L); \node[vlab,left]  at (l2L) {$r_2'$};
  \draw[fig] (O) -- (kb);  \node[vlab,right] at (kb) {$l_1$};
  \draw[aux] (O) -- ($(O)!1.18!(ka)$);
  \node[vlab,below] at ($(O)!1.18!(ka)$) {$r_1'$};
  \node[vlab,below] at ($(O)!0.62!(l2R)$) {$r_2$};
  \node[vlab] at ($($(O)!0.78!(kb)$)+(147.8:0.42)$) {$r_1$};
  % Q_1 is on the same side of l_2 as r_1, the far side of l_1 from r_2;
  % R is where segment Q_1 P meets l_1 -- it has to land on the ray r_1.
  \coordinate (P) at (4.05,2.10); \coordinate (Q) at (1.35,2.45);
  \draw[fig] (Q) -- (P);
  \coordinate (R) at (intersection of Q--P and ka--kb);
  \foreach \p in {P,Q,R} {\dt{\p}}
  \node[vlab,right] at (P) {$P$}; \node[vlab,left] at (Q) {$Q_1$};
  \node[vlab] at ($(R)+(299.5:0.55)$) {$R$};
\end{tikzpicture}
\figcap{3--38}
\end{minipage}\hfill
\begin{minipage}{0.48\linewidth}\centering
\begin{tikzpicture}[scale=0.675]
  \coordinate (l2L) at (0.30,1.05); \coordinate (l2R) at (4.55,1.25);
  \coordinate (ka) at (1.65,0.15);  \coordinate (kb) at (3.35,2.85);
  \coordinate (O) at (intersection of l2L--l2R and ka--kb);
  \draw[fig] (O) -- (l2R); \node[vlab,right] at (l2R) {$l_2$};
  \draw[aux] (O) -- (l2L); \node[vlab,left]  at (l2L) {$r_2'$};
  \draw[fig] (O) -- (kb);  \node[vlab,right] at (kb) {$l_1$};
  \draw[aux] (O) -- ($(O)!1.18!(ka)$);
  \node[vlab,below] at ($(O)!1.18!(ka)$) {$r_1'$};
  \node[vlab,below] at ($(O)!0.62!(l2R)$) {$r_2$};
  \node[vlab] at ($($(O)!0.78!(kb)$)+(147.8:0.42)$) {$r_1$};
  % the one thing that separates this figure from 3-38: Q_2 lies ON the ray
  % r_2' itself, not merely on that side of it
  \coordinate (P) at (4.05,2.10); \coordinate (Q) at ($(l2L)!0.16!(l2R)$);
  \draw[fig] (Q) -- (P);
  \coordinate (R) at (intersection of Q--P and ka--kb);
  \foreach \p in {P,Q,R} {\dt{\p}}
  \node[vlab,right] at (P) {$P$}; \node[vlab,above] at (Q) {$Q_2$};
  \node[vlab] at ($(R)+(126:0.60)$) {$R$};
\end{tikzpicture}
\figcap{3--39}
\end{minipage}
\end{bkfigure}}
